\documentclass[11pt,a4paper]{article}
\usepackage[utf8]{inputenc}
\usepackage[T1]{fontenc}
\usepackage[portuguese]{babel}
\usepackage{xcolor}
\usepackage{hyperref}
\usepackage{geometry}
\usepackage{tcolorbox}

\geometry{margin=2.5cm}

\definecolor{primary}{RGB}{41, 128, 185}
\definecolor{secondary}{RGB}{44, 62, 80}
\definecolor{codebg}{RGB}{240, 240, 240}

\hypersetup{colorlinks=true, linkcolor=primary, urlcolor=primary}

\title{\color{secondary}\Huge\textbf{Exercício 4.2: O projeto, passo 21}}
\author{\color{primary}DevOps com Kubernetes -- 2026}
\date{}

\begin{document}
\maketitle

\section*{\color{primary}Instruções do Exercício}
\begin{tcolorbox}[colback=blue!5!white,colframe=blue!75!black,title=Exercício 4.2: O projeto, passo 21]
Crie as sondas de prontidão e vivacidade (\texttt{ReadinessProbes} e \texttt{LivenessProbes}), bem como o \textit{endpoint} necessário (como \texttt{/healthz}) para O Projeto, de forma a garantir que a aplicação esteja funcionando corretamente e conectada ao banco de dados.

Adicione um botão à sua interface que pode ser usado para 'quebrar' a aplicação (simular uma falha). Pressionar o botão deve fazer com que a operação normal pare.

Certifique-se de que, uma vez quebrado, o Kubernetes identifique a falha e um novo pod inicie, ficando saudável novamente.
\end{tcolorbox}

\end{document}
